% Copyright 2020 Palle Jorgensen, hamselv@pallej.dk
%
% This file is part of utf8add.
%
% utf8add is free software: you can redistribute it and/or modify it
% under the terms of the GNU General Public License as published by
% the Free Software Foundation, either version 3 of the License, or
% (at your option) any later version.
%
% utf8add is distributed in the hope that it will be useful, but WITHOUT
% ANY WARRANTY; without even the implied warranty of MERCHANTABILITY
% or FITNESS FOR A PARTICULAR PURPOSE. See the GNU General Public
% License for more details.
%
% You should have received a copy of the GNU General Public License
% along with utf8add. If not, see <https://www.gnu.org/licenses/>.
\documentclass[a4paper,british,parskip=half,11pt,fleqn]{scrartcl}
\usepackage[T1]{fontenc} %
\usepackage[utf8]{inputenc} %
\usepackage{babel} %
\usepackage[sc]{mathpazo} %
\usepackage{tgpagella} %
\usepackage{listings} %
\usepackage[scaled]{luximono} %
\usepackage{textcomp} %
\usepackage{classico} %
\linespread{1.045} %
\usepackage[final]{microtype} %
%\setcounter{secnumdepth}{1} %
\author{Palle J\o rgensen} %
\title{The \texttt{\jobname} package} %
\subtitle{Using UTF-8 encoded input as short cuts} %
\lstset{language=[latex]tex}
\setkomafont{descriptionlabel}{\ttfamily\bfseries}
\let\oldquote\quote
\def\quote{\leftmargini0pt\oldquote\slshape}

\begin{document}
\maketitle

\section{Introduction}
\label{sec:introduction}

The purpose of this package is to provide som short cuts for typing
math in \LaTeX. It is \emph{not} intended for increasing readability.

\section{Loading the package}
\label{sec:loading-package}

\begin{lstlisting}
\usepackage[default]{utf8hax}
\end{lstlisting}

\subsection{Package options}
\label{sec:package-options}

\begin{description}
\item[autochar] Choosing the character for the math automatation. It
  must be the hex code for the desired character.
\item[mathchar] Choosing the character for math mode. It must be the
  hex code for the desired character.
\item[currmath] Setting `\verb|¤|' as the math character.
\item[yenmath]  Setting `\verb|¥|' as the math character.
\item[sterlingmath] Setting `\verb|£|' as the math character.
\item[copyauto] Setting `\verb|©|' at the math automatation character.
\item[regauto] Setting `\verb|®|' at the math automatation character.
\item[default] Short cut for the \verb|copyauto| and \verb|currmath|
  option. The authors favourite.
\end{description}

\section{Using the package}
\label{sec:using-package}

\subsection{Math character}
\label{sec:math-character}

Typing 
\begin{verbatim}
The slope of the line is designated ¤a.
\end{verbatim}
equals
\begin{verbatim}
The slope of the line is designated \ensuremath{a}.
\end{verbatim}
yielding
\begin{quote}
  The slope of the line is designated \ensuremath{a}.
\end{quote}

\subsection{Math automatation character}
\label{sec:math-autom-char}

Typing 
\begin{verbatim}
The slope of the line is designated ©a.
\end{verbatim}
automates into
\begin{verbatim}
The slope of the line is designated \@automath@a.
\end{verbatim}
The command `\verb|\@automath@a|' is programmed into being equal to `\verb|\ensuremath{\alpha}|'.

Thus the input is equal to
\begin{verbatim}
The slope of the line is designated \ensuremath{\alpha}.
\end{verbatim}
yielding
\begin{quote}
The slope of the line is designated \ensuremath{\alpha}.
\end{quote}

\subsection{New math automatation}
\label{sec:new-math-shortcuts}

New math automatation is accessible through the
`\verb|\newmathautomatation|' command. Typing
\begin{verbatim}
\newmathautomatation{?}{What?}
\end{verbatim}
makes the input sequence `\verb|©?|' equal to
\begin{verbatim}
\ensuremath{What?}
\end{verbatim}
yielding
\begin{quote}
\ensuremath{What?}
\end{quote}

\section{Drawbacks}
\label{sec:drawbacks}

As some input characters are changed the access to these characters is
made more difficult.

The copyright symbol is accessible through `\verb|©!|'. As i do not
know anybody using the \emph{generic currency symbol} in real life it
is not made as easily accessible.

\section{Copyright and license}
\label{sec:copyright-license}

This file \texttt{\jobname} is part of the \texttt{utf8add} bundle.

The \texttt{utf8add} package and all the associated files are
copyright © Palle Jørgensen, \ifnum\year=2020\relax 2020\else
2020-\the\year\fi.

\texttt{utf8add} is free software: you can redistribute it and/or
modify it under the terms of the GNU General Public License as
published by the Free Software Foundation, either version 3 of the
License, or (at your option) any later version.

\texttt{utf8add} is distributed in the hope that it will be useful,
but WITHOUT ANY WARRANTY; without even the implied warranty of
MERCHANTABILITY or FITNESS FOR A PARTICULAR PURPOSE. See the GNU
General Public License for more details.

You should have received a copy of the GNU General Public License
along with \texttt{utf8add}. If not, see
<https://www.gnu.org/licenses/>.


\end{document}

%%% Local Variables:
%%% mode: latex
%%% TeX-master: t
%%% End:
